\chapter{Natural Numbers and Processing Multiple Lists}

\section{Natural numbers as recursive data}

Lists are built from \code{empty} using \code{cons}. Similarly, natural numbers
can be viewed as built from \code{0} using \code{add1}:

\begin{lstlisting}[style=racket]
;; A Nat is one of:
;; * 0
;; * (add1 Nat)
\end{lstlisting}

The analogues of \code{empty?} and \code{rest} are \code{zero?} and
\code{sub1}. This gives a template:

\begin{lstlisting}[style=racket]
;; nat-template: Nat -> Any
(define (nat-template n)
  (cond [(zero? n) ...]
        [else (... n ... (nat-template (sub1 n)) ...)]))
\end{lstlisting}

\begin{example}[Countdown]
\begin{lstlisting}[style=racket]
;; (countdown n) produces the list n, n-1, ..., 1, 0.
;; Examples:
(check-expect (countdown 0) '(0))
(check-expect (countdown 3) '(3 2 1 0))

;; countdown: Nat -> (listof Nat)
(define (countdown n)
  (cond [(zero? n) (cons 0 empty)]
        [else (cons n (countdown (sub1 n)))]))
\end{lstlisting}
\end{example}

\begin{example}[Factorial]
\begin{lstlisting}[style=racket]
;; (factorial n) produces n! = 1 * 2 * ... * n.
;; factorial: Nat -> Nat
(check-expect (factorial 0) 1)
(check-expect (factorial 5) 120)
(define (factorial n)
  (cond [(zero? n) 1]
        [else (* n (factorial (sub1 n)))]))
\end{lstlisting}
Condensed trace:
\begin{lstlisting}[style=racket]
(factorial 3)
=> (* 3 (factorial 2))
=> (* 3 (* 2 (factorial 1)))
=> (* 3 (* 2 (* 1 (factorial 0))))
=> (* 3 (* 2 (* 1 1)))
=> 6
\end{lstlisting}
\end{example}

\subsection{Subintervals of the naturals}
To count down to a lower bound \code{b} instead of 0, change the base case:
\begin{lstlisting}[style=racket]
;; An integer n >= b is one of:
;; * b
;; * (add1 m) where m >= b

;; (countdown-to n b) produces n, n-1, ..., b.
;; countdown-to: Int Int -> (listof Int)
;;   requires: n >= b
(check-expect (countdown-to 5 3) '(5 4 3))
(check-expect (countdown-to -1 -2) '(-1 -2))
(define (countdown-to n b)
  (cond [(= n b) (cons b empty)]
        [else (cons n (countdown-to (sub1 n) b))]))
\end{lstlisting}

Counting \emph{up} uses \code{add1} toward an upper bound:
\begin{lstlisting}[style=racket]
;; (countup-to n b) produces n, n+1, ..., b.
;; countup-to: Int Int -> (listof Int)
;;   requires: n <= b
(check-expect (countup-to 2 5) '(2 3 4 5))
(define (countup-to n b)
  (cond [(= n b) (cons b empty)]
        [else (cons n (countup-to (add1 n) b))]))
\end{lstlisting}

\begin{example}[Primality by trial division]
\begin{lstlisting}[style=racket]
;; (prime? n) produces true if n is prime.
;; prime?: Nat -> Bool
(check-expect (prime? 7) true)
(check-expect (prime? 9) false)
(check-expect (prime? 1) false)
(define (prime? n)
  (and (> n 1) (no-divisors? n (sub1 n))))

;; (no-divisors? n k) produces true if no integer in 2..k divides n.
;; no-divisors?: Nat Nat -> Bool
(define (no-divisors? n k)
  (cond [(<= k 1) true]
        [else (and (not (zero? (remainder n k)))
                   (no-divisors? n (sub1 k)))]))
\end{lstlisting}
\end{example}

\section{Processing two lists}

When a function consumes two lists, there are three typical situations.

\subsection{Case 1: one list goes along for the ride}
Only one list is recursed on; the other is passed unchanged. \code{my-append}
from the previous chapter is the standard example. The template is the
ordinary list template with an extra parameter.

\subsection{Case 2: processing in lockstep}
The lists have the same length, and corresponding elements are processed
together.
\begin{lstlisting}[style=racket]
;; (dot-product l1 l2) produces the dot product of l1 and l2.
;; Examples:
(check-expect (dot-product '(1 2 3) '(4 5 6)) 32)
(check-expect (dot-product empty empty) 0)

;; dot-product: (listof Num) (listof Num) -> Num
;;   requires: l1 and l2 have the same length
(define (dot-product l1 l2)
  (cond [(empty? l1) 0]
        [else (+ (* (first l1) (first l2))
                 (dot-product (rest l1) (rest l2)))]))
\end{lstlisting}
Since both lists empty out simultaneously, checking one suffices.

\subsection{Case 3: processing at different rates}
In general, either list can be empty independently, giving four combinations.

\begin{lstlisting}[style=racket]
;; twolist-template: (listof X) (listof Y) -> Any
(define (twolist-template lx ly)
  (cond [(and (empty? lx) (empty? ly)) ...]
        [(and (empty? lx) (cons? ly))
         (... (first ly) ... (rest ly) ...)]
        [(and (cons? lx) (empty? ly))
         (... (first lx) ... (rest lx) ...)]
        [(and (cons? lx) (cons? ly))
         ;; may recurse on (rest lx), (rest ly), or both
         (... (first lx) ... (first ly) ...
              (twolist-template (rest lx) ly) ...
              (twolist-template lx (rest ly)) ...
              (twolist-template (rest lx) (rest ly)) ...)]))
\end{lstlisting}
Often some cases collapse.

\begin{example}[Merging sorted lists]
\begin{lstlisting}[style=racket]
;; (merge l1 l2) produces a sorted list of all elements of the sorted
;;   lists l1 and l2.
;; Examples:
(check-expect (merge '(1 4 9) '(2 3 10)) '(1 2 3 4 9 10))
(check-expect (merge empty '(5)) '(5))

;; merge: (listof Num) (listof Num) -> (listof Num)
;;   requires: l1 and l2 are sorted in non-decreasing order
(define (merge l1 l2)
  (cond [(empty? l1) l2]
        [(empty? l2) l1]
        [(< (first l1) (first l2))
         (cons (first l1) (merge (rest l1) l2))]
        [else (cons (first l2) (merge l1 (rest l2)))]))
\end{lstlisting}
Condensed trace:
\begin{lstlisting}[style=racket]
(merge (list 1 4) (list 2 3))
=> (cons 1 (merge (list 4) (list 2 3)))
=> (cons 1 (cons 2 (merge (list 4) (list 3))))
=> (cons 1 (cons 2 (cons 3 (merge (list 4) empty))))
=> (cons 1 (cons 2 (cons 3 (list 4))))
=> (list 1 2 3 4)
\end{lstlisting}
\end{example}

\begin{example}[List equality]
\begin{lstlisting}[style=racket]
;; (list=? l1 l2) produces true if l1 and l2 contain the same numbers
;;   in the same order.
;; list=?: (listof Num) (listof Num) -> Bool
(check-expect (list=? '(1 2) '(1 2)) true)
(check-expect (list=? '(1 2) '(1)) false)
(define (list=? l1 l2)
  (cond [(and (empty? l1) (empty? l2)) true]
        [(or (empty? l1) (empty? l2)) false]
        [else (and (= (first l1) (first l2))
                   (list=? (rest l1) (rest l2)))]))
\end{lstlisting}
\end{example}

\subsection{Consuming a list and a number}
A list and a \code{Nat} can be processed together, e.g.\ \code{list-ref}:
\begin{lstlisting}[style=racket]
;; (my-list-ref lst n) produces the element at position n (0-based),
;;   or false if lst is too short.
;; my-list-ref: (listof Any) Nat -> (anyof Any false)
(check-expect (my-list-ref '(a b c) 1) 'b)
(check-expect (my-list-ref '(a) 3) false)
(define (my-list-ref lst n)
  (cond [(empty? lst) false]
        [(zero? n) (first lst)]
        [else (my-list-ref (rest lst) (sub1 n))]))
\end{lstlisting}

\section{Summary}
\begin{itemize}
  \item Natural numbers have a recursive data definition (\code{0} or \code{add1}); template uses \code{zero?} and \code{sub1}.
  \item Intervals: change the base case to a bound \code{b}; counting up uses \code{add1}.
  \item Two lists: along for the ride, lockstep, or at different rates (four-case template).
  \item \code{merge} of sorted lists is the key example of different-rate processing.
\end{itemize}
